\documentclass[12pt,a4paper]{article}
\usepackage{amsmath,amssymb,amsthm}
\usepackage{mathtools}
\usepackage{hyperref}
\usepackage{geometry}
\usepackage{booktabs}
\usepackage{enumitem}
\usepackage{xcolor}
\geometry{margin=2.5cm}
\hypersetup{colorlinks=true,linkcolor=blue,citecolor=blue,urlcolor=blue}

\title{%
  SOLAR/DUNE Significance Analyses\\
  \large Mathematical Derivations: Day-Night, HEP, and Sensitivity%
}
\author{SOLAR/DUNE Analysis}
\date{May 2026}

\newcommand{\bhat}{\hat{\beta}}
\newcommand{\srel}{\sigma_{\mathrm{rel}}}
\newcommand{\srelsq}{\sigma_{\mathrm{rel}}^{2}}
\newcommand{\Btot}{B_{\mathrm{tot}}}
\newcommand{\Ntot}{N_{\mathrm{tot}}}
\newcommand{\minexp}{\epsilon_{\min}}
\newcommand{\Tsig}{T^{\mathrm{sig}}}
\newcommand{\Tbkg}{T^{\mathrm{bkg}}}
\newcommand{\Apred}{A_{\mathrm{pred}}}
\newcommand{\Abkg}{A_{\mathrm{bkg}}}
\newcommand{\spred}{\sigma_{\mathrm{pred}}}
\newcommand{\sbkg}{\sigma_{\mathrm{bkg}}}

\begin{document}
\maketitle
\tableofcontents
\newpage

% ===========================================================================
\section{Overview and Scientific Goals}
\label{sec:overview}

This document presents the mathematical derivations underlying the three
significance analyses in the SOLAR/DUNE software framework,
ordered from simplest to most complex:

\begin{enumerate}[leftmargin=*]
  \item \textbf{Day-Night Asymmetry} (\texttt{12DayNight.py}):
    searches for a time-modulated excess of solar neutrinos at night
    relative to the day.  The significance is computed using a
    Gaussian approximation on the difference spectrum.

  \item \textbf{HEP Discovery} (\texttt{13HEP.py}):
    searches for the hep solar neutrino flux ($^3\mathrm{He}+p \to
    {}^4\mathrm{He}+e^++\nu_e$) as an absolute excess above background.
    The primary metric is the profile-likelihood (PL) significance
    with a single global background nuisance, analytically profiled.

  \item \textbf{Sensitivity} (\texttt{14Sensitivity.py}):
    maps the 2D sensitivity contour in the oscillation-parameter plane
    $(\Delta m^2,\,\sin^2\theta_{12})$.
    A Baker-Cousins Poisson deviance with two normalization nuisances
    is minimized numerically over 2D templates in energy and azimuth.
\end{enumerate}

The three analyses share common ingredients (smoothing, thresholds,
adaptive rebinning) but differ in the complexity of their statistical model.
Table~\ref{tab:comparison} in Section~\ref{sec:comparison} summarizes
the key differences.

\medskip
\noindent\textbf{Common notation.}
Throughout, $T$ denotes the analysis exposure in kt$\cdot$yr,
$M_{\mathrm{det}}$ the active detector mass in kt,
and the exposure factor is $\mathcal{E} = T\,M_{\mathrm{det}}$.
Index $i$ runs over energy bins, index $j$ over azimuth bins (Sensitivity only).

% ===========================================================================
\section{Common Ingredients}
\label{sec:common}

\subsection{Rate Histograms}
\label{sec:rates}

All three analyses work with per-unit-exposure, per-unit-mass rate histograms
$r_i^c$ (units: $(\mathrm{yr\cdot kt\cdot MeV})^{-1}$, integrated over bin width)
for each component $c$ (signal, neutrons, gammas, radiologicals, solar $^8$B).
Event counts are obtained by scaling with the exposure factor:
\begin{equation}
  \mu_i^c = \mathcal{E}\, r_i^c = T\,M_{\mathrm{det}}\, r_i^c.
  \label{eq:rate_to_counts}
\end{equation}
All analyses apply a reconstructed-energy threshold $E_{\mathrm{th}}$
(configured in \texttt{import/analysis.json}); only bins with $E_i \ge E_{\mathrm{th}}$
enter the significance computation.

\subsection{Histogram Smoothing}
\label{sec:smoothing}

Each component histogram $r_i^c$ is convolved with a one-dimensional Gaussian kernel
before entering the significance computation:
\begin{equation}
  \tilde{r}_{i}^{c} = \sum_{j} G_\sigma(i-j)\, r_j^{c},
  \qquad
  G_\sigma(k) = \frac{1}{\sqrt{2\pi}\,\sigma}
    \exp\!\left(-\frac{k^2}{2\sigma^2}\right),
  \label{eq:gaussian_smooth}
\end{equation}
implemented via \texttt{scipy.ndimage.gaussian\_filter1d} with
\texttt{mode=`nearest'}.  The smoothing width $\sigma$ and whether
smoothing is applied per component are configured in
\texttt{import/analysis.json} under \texttt{SMOOTHING.ANALYSES.\{ANALYSIS\}.STAGES}.

\medskip
\noindent\textbf{Rate clipping.}
Gaussian convolution at distribution tails can produce small negative values.
Negative rates are unphysical and cause profile-likelihood divergences
at high exposure (Section~\ref{sec:pl_numerical}).
All smoothed rates are therefore clipped:
\begin{equation}
  \tilde{r}_{i}^{c} \leftarrow \max\!\left(0,\,\tilde{r}_{i}^{c}\right).
  \label{eq:clip}
\end{equation}

\subsection{Background Error Propagation}
\label{sec:bkg_error}

The absolute uncertainty on the raw background rate in bin $i$ from component $c$
combines MC-statistical and systematic contributions:
\begin{equation}
  \sigma_{i}^{c} = r_i^{c}\,
    \sqrt{\left(\frac{\delta_i^c}{r_i^c}\right)^{\!2}
          + \left(\srel^c\right)^2},
  \label{eq:bin_error}
\end{equation}
where $\delta_i^c$ is the histogram statistical error and $\srel^c$ is
the relative systematic uncertainty (\texttt{--background\_uncertainty}).
The combined background uncertainty is $\sigma_i = \sqrt{\sum_c (\sigma_i^c)^2}$,
also smoothed with the component kernel.

% ===========================================================================
\section{Day-Night Asymmetry Analysis}
\label{sec:daynight}

\subsection{Physical Observable: MSW-Enhanced Day-Night Effect}
\label{sec:dn_physics}

Solar neutrinos produced as $\nu_e$ oscillate in vacuum en route to the Earth.
During nighttime passages, neutrinos cross the Earth's interior and experience
additional matter-induced (MSW) oscillations that partially restore the
$\nu_e$ component.  This creates a nighttime excess:
\begin{equation}
  \Delta_{\mathrm{DN}} \equiv \frac{N_{\mathrm{night}} - N_{\mathrm{day}}}
    {{\tfrac{1}{2}(N_{\mathrm{night}} + N_{\mathrm{day}})}},
  \label{eq:dn_asymmetry}
\end{equation}
where $N_{\mathrm{day/night}}$ are the total event counts in the
corresponding half-runs.  The size of $\Delta_{\mathrm{DN}}$ depends
on the Earth density profile through the MSW matter-effect potential
$V = \sqrt{2}\,G_F\, n_e(r)$, and on the oscillation parameters
$(\Delta m^2,\,\sin^2\theta_{12})$.

\subsection{Day-Night Signal and Background Model}
\label{sec:dn_model}

The analysis works per-energy-bin with the oscillation-weighted solar
signal split into day and night components:
\begin{align}
  r_i^{\mathrm{day}} &= \text{solar rate, daytime sky exposure},\\
  r_i^{\mathrm{night}} &= \text{solar rate, nighttime (Earth-crossing) exposure}.
\end{align}
Isotropic backgrounds (cosmogenic, geological, detector-intrinsic) are assumed
time-uniform; they split equally between day and night runs
($f_{\mathrm{day}} = 0.5$).

The \textbf{asymmetry signal} per bin at exposure $\mathcal{E}$ is:
\begin{equation}
  s_i = \mathcal{E}\cdot\xi\cdot\bigl(r_i^{\mathrm{night}} - r_i^{\mathrm{day}}\bigr),
  \label{eq:dn_signal}
\end{equation}
where $\xi$ is the Earth density scale factor (Section~\ref{sec:dn_band}).
The \textbf{effective background} for the asymmetry measurement is:
\begin{equation}
  b_i = \mathcal{E}\cdot\bigl(f_{\mathrm{day}}\,r_i^{\mathrm{bkg}} + r_i^{\mathrm{day}}\bigr),
  \label{eq:dn_background}
\end{equation}
where the daytime solar rate $r_i^{\mathrm{day}}$ acts as background
because it is present symmetrically and cannot be distinguished from
day-independent sources in the asymmetry.  Bins where $b_i=0$ are zeroed.

\subsection{Gaussian Significance}
\label{sec:dn_gaussian}

The Day-Night analysis uses Gaussian significance only; no profile-likelihood
is applied.  Two variants are computed:

\medskip
\noindent\textbf{Simple Gaussian} (systematic uncertainty only):
\begin{equation}
  Z_i^{\mathrm{G}} = \frac{s_i}{\sqrt{b_i}},
  \label{eq:dn_simple_gaussian}
\end{equation}

\noindent\textbf{Error Gaussian} (systematic + Poisson statistical uncertainty):
\begin{equation}
  Z_i^{\mathrm{G,err}} = \frac{s_i}{\sqrt{b_i + \sigma_i^2}},
  \qquad
  \sigma_i = \sqrt{\mathcal{E}\,f_{\mathrm{day}}\,r_i^{\mathrm{bkg}}},
  \label{eq:dn_error_gaussian}
\end{equation}
where $\sigma_i$ is the absolute Poisson uncertainty on the isotropic
background count in the daytime half-run.  The global significance is the
quadrature sum:
\begin{equation}
  Z = \sqrt{\sum_{i \ge i_{\mathrm{th}}} \!\left(Z_i\right)^2}.
  \label{eq:dn_global}
\end{equation}
Both raw and Gaussian-smoothed versions of the rates are used,
producing four output curves per cut
(\texttt{RawGaussian}, \texttt{RawErrorGaussian}, \texttt{Gaussian},
\texttt{ErrorGaussian}) plus $\pm$Earth-density-band variants.

\subsection{Earth Density Uncertainty Bands}
\label{sec:dn_band}

The MSW matter effect depends on the Earth's electron density profile.
Published PREM-based oscillation probability calculations span a
$\pm\varepsilon$ range around the nominal prediction.
Three values of the scale factor $\xi$ are used:
\begin{equation}
  \xi \in \{1+\varepsilon,\;1,\;1-\varepsilon\},
  \qquad
  \varepsilon = 0.13\;\text{(default, }\pm 13\%\text{)}.
  \label{eq:dn_band}
\end{equation}
The scale acts on the \emph{asymmetry} only (Eq.~\eqref{eq:dn_signal}),
leaving the total normalization and the background unchanged.
This produces three significance curves: upper, nominal, lower Earth-density
predictions, bracketing the true uncertainty on the matter-effect amplitude.

% ===========================================================================
\section{HEP Discovery Analysis}
\label{sec:hep}

\subsection{Signal Model and Exposure Scaling}
\label{sec:hep_signal}

The HEP analysis searches for an absolute rate excess of hep solar neutrinos
above a known background in a one-dimensional energy spectrum above $E_{\mathrm{th}}$.
Signal and background event counts in bin $i$ at exposure $\mathcal{E}$ are:
\begin{align}
  s_i &= \mathcal{E}\, r_i^{\mathrm{sig}},
  \label{eq:hep_signal}\\
  b_i &= \mathcal{E}\, r_i^{\mathrm{bkg}}.
  \label{eq:hep_background}
\end{align}
The Asimov dataset sets observed counts to $n_i = s_i + b_i$.

\subsection{Adaptive Tail Rebinning}
\label{sec:rebin}

\subsubsection{Detection Threshold}
\label{sec:detection_threshold}

A bin (or merged group) is declared \emph{detectable} if the expected
signal exceeds:
\begin{equation}
  T_{\mathrm{det}} = \max\!\left(T_{\mathrm{min}},\;-\ln(1-p_{\mathrm{min}})\right),
  \label{eq:detection_threshold}
\end{equation}
with defaults $T_{\mathrm{min}}=1.0$ and $p_{\mathrm{min}}=1-e^{-1}$,
giving $T_{\mathrm{det}}=1.0$.  The Poisson interpretation:
$P(\ge 1;\lambda)\ge p \Leftrightarrow \lambda \ge -\ln(1-p)$.
The conservative detectable signal uses a downward fluctuation:
\begin{equation}
  s_i^{\mathrm{det}} = s_i(1 - d\cdot\srel^{\mathrm{sig}}),
  \quad d \in \{2.9,3.0,3.1\}.
  \label{eq:detectable_signal}
\end{equation}

\subsubsection{Greedy Algorithm and Monotonicity Enforcement}

Bins are merged from the high-energy tail inward until each group exceeds
$T_{\mathrm{det}}$.  Groups below $T_{\mathrm{det}}$ are zeroed.
To guarantee a non-decreasing significance vs exposure curve,
if the current step's optimal significance is below the previous:
\begin{equation}
  \mathcal{B}(t) =
  \begin{cases}
    \mathcal{B}^*(t) & Z_A(\mathcal{B}^*(t),t) \ge Z_A(\mathcal{B}(t-1),t-1),\\
    \mathcal{B}(t-1) & \text{otherwise.}
  \end{cases}
  \label{eq:monotone_rebin}
\end{equation}

\subsection{Gaussian and Asimov Significance}
\label{sec:hep_gaussian_asimov}

The per-bin Gaussian significance and its quadrature sum follow
Eqs.~\eqref{eq:dn_simple_gaussian}--\eqref{eq:dn_global}, applied to the
absolute signal (not a difference spectrum).
The Asimov significance \cite{Cowan2010}, with and without background uncertainty, is:
\begin{align}
  Z_A &= \sqrt{2\left[(s+b)\ln\!\left(1+\tfrac{s}{b}\right) - s\right]},
  \label{eq:asimov_simple}\\
  Z_A^{\sigma} &= \sqrt{2\left[
    (s+b)\ln\!\frac{(s+b)(b+\sigma^2)}{b^2+(s+b)\sigma^2}
    -\frac{b^2}{\sigma^2}\ln\!\left(1+\frac{\sigma^2 s}{b(b+\sigma^2)}\right)
  \right]},
  \label{eq:asimov_full}
\end{align}
both applied per-bin with the global significance as the quadrature sum.
These are used for cross-checks; the primary metric is the PL significance below.

\subsection{Profile-Likelihood Significance}
\label{sec:pl}

\subsubsection{Test Statistic}
\label{sec:pl_model}

The profile-likelihood significance follows Cowan et al.\ \cite{Cowan2010}.
For the null hypothesis $\mu=0$ (no HEP signal), the test statistic is:
\begin{equation}
  q_0 = -2\ln\!\frac{\mathcal{L}(0,\,\bhat)}{\mathcal{L}(\hat{s}+b,\,1)},
  \label{eq:q0}
\end{equation}
evaluated at the Asimov point $n_i = s_i + b_i$, with $\bhat$ the
background nuisance profiled under the null.
The median discovery significance is $Z = \sqrt{q_0}$.

\subsubsection{Background Normalization Nuisance}
\label{sec:pl_nuisance}

A single global background scale factor $\beta$ is constrained by a
Gaussian penalty:
\begin{equation}
  \mathcal{L}(0,\beta)
  = \prod_{i=1}^{K}
      \frac{(\beta b_i)^{n_i} e^{-\beta b_i}}{n_i!}
    \cdot\exp\!\left(-\frac{(\beta-1)^2}{2\srelsq}\right),
  \label{eq:likelihood}
\end{equation}
where $\srel$ is the relative background uncertainty
(\texttt{--background\_uncertainty}, typically 2\%).
A global (fully correlated) $\beta$ avoids the $K$-parameter absorbing
plateau that per-bin nuisances produce; the plateau ends at
$T \sim 1/(B_{\mathrm{tot}}\srelsq)$ (typically sub-year), giving
a physically correct PL curve.

\subsubsection{Solving for \texorpdfstring{$\bhat$}{beta-hat}}
\label{sec:betahat}

The stationarity condition $\partial\ln\mathcal{L}/\partial\beta|_{\bhat}=0$
at the Asimov point $n_i=s_i+b_i$ gives:
\begin{equation}
  \Ntot = \bhat\Btot + \frac{\bhat-1}{\srelsq},
  \quad
  \Ntot = \sum_i n_i,\;\Btot = \sum_i b_i.
  \label{eq:stationarity}
\end{equation}
This is the key feature distinguishing HEP from Sensitivity: because $\beta$
is \emph{global}, all per-bin denominators factor as $\beta b_i$, and the
stationarity condition collapses to a single equation in one unknown.
Rearranging gives a quadratic:
\begin{equation}
  \bhat^2 + \left(\Btot\srelsq-1\right)\bhat - \Ntot\srelsq = 0,
  \label{eq:beta_quadratic}
\end{equation}
with analytic positive root:
\begin{equation}
  \bhat = \frac{-(B_{\mathrm{tot}}\srelsq-1)
    + \sqrt{(B_{\mathrm{tot}}\srelsq-1)^2 + 4\Ntot\srelsq}}{2}.
  \label{eq:beta_solution}
\end{equation}
When $\Btot=0$ or $\srel=0$, $\bhat=1$ (no constraint active).

\subsubsection{Log-Likelihood Ratio}
\label{sec:llr}

Substituting $n_i = s_i+b_i$ into Eq.~\eqref{eq:q0}:
\begin{equation}
  q_0 = 2\sum_{i=1}^{K}
    \left[n_i\ln\!\frac{n_i}{\bhat b_i} - (n_i-\bhat b_i)\right]
    + \left(\frac{\bhat-1}{\srel}\right)^{\!2}.
  \label{eq:q0_expanded}
\end{equation}
Each per-bin term is the Baker-Cousins Poisson deviance~\cite{BakerCousins1984}
between observed $n_i$ and null expectation $\bhat b_i$.
This connection — the HEP LLR is a sum of per-bin Poisson deviances with a
global nuisance penalty — foreshadows the Sensitivity objective function
(Section~\ref{sec:sens_objective}).

\subsubsection{Numerical Stability}
\label{sec:pl_numerical}

Naive subtraction of two $O(N\log N)$ log-likelihoods loses precision when
$S^2/B\ll 1$.  The stable per-bin form:
\begin{equation}
  \Delta\ell_i = n_i\ln\!\frac{n_i}{\mu_i^{\mathrm{null}}}
    - \bigl(n_i - \mu_i^{\mathrm{null}}\bigr),
  \qquad
  \mu_i^{\mathrm{null}} = \bhat b_i,
  \label{eq:stable_llr}
\end{equation}
with a floor $\minexp=10^{-12}$ replacing zero denominators.

\medskip
\noindent\textbf{Negative-rate blowup.}
If a smoothed rate becomes negative, $\bhat b_i<0$, the floor gives
$\Delta\ell_i \approx n_i\ln(n_i/\minexp) \propto T$,
causing a spurious superlinear spike.
Clipping (Eq.~\eqref{eq:clip}) eliminates this.

\subsection{Barlow-Beeston MC Mask}
\label{sec:bb_mask}

Bins with summed background MC count below $N_{\mathrm{MC}}^{\min}$ (default 1.0)
lack a reliable background model and are excluded by a static binary mask
\cite{BarlowBeeston1993}:
\begin{equation}
  \mathcal{M}_i = \mathbf{1}\!\left[N_i^{\mathrm{MC}} \ge N_{\mathrm{MC}}^{\min}\right],
  \qquad
  s_i \leftarrow \mathcal{M}_i s_i,\quad b_i \leftarrow \mathcal{M}_i b_i.
  \label{eq:bb_mask_hep}
\end{equation}
The static mask introduces no discrete transitions in the exposure curve.

\subsection{Profile-Likelihood Error Bands}
\label{sec:pl_bands}

The $\pm1\sigma$ bands use signal normalization variation with
$\delta\in\{+1,0,-1\}$:
\begin{equation}
  s_i^{(\delta)} = s_i(1 + \delta\,\srel^{\mathrm{sig}}).
  \label{eq:signal_variation}
\end{equation}
Background is never shifted, so $\bhat$ is unaffected and both bands
collapse symmetrically when the signal is negligible.

\subsection{Post-Processing: PL Smoothing and Monotonicity}
\label{sec:pl_smooth}

The raw PL significance curve $Z(T)$ is post-processed in two steps:
\begin{enumerate}[leftmargin=*]
  \item \textbf{Gaussian kernel smoothing} with $\sigma_{\mathrm{PL}}=6$
    exposure-grid index units (\texttt{scipy.ndimage.gaussian\_filter1d},
    equivalent to ROOT's \texttt{TH1::Smooth} on the exposure axis).
  \item \textbf{Isotonic regression (PAVA)} \cite{Robertson1988}:
    \begin{equation}
      \min_{\{z_k\}}\sum_k(z_k-\tilde{Z}(T_k))^2
      \quad\text{s.t.}\quad z_k\le z_{k+1},
      \label{eq:pava}
    \end{equation}
    via \texttt{sklearn.isotonic.IsotonicRegression}.
\end{enumerate}
Pipeline: $Z \to \max(0,Z) \xrightarrow{\mathrm{Gauss}} \tilde{Z} \xrightarrow{\mathrm{PAVA}} Z'$.

\subsection{Spike Detection and Best-Cut Selection}
\label{sec:spike}

A PL curve is \emph{spiked} if any consecutive step in the pre-isotonic columns
(\texttt{RawPreIsotonicProfileLikelihood} or \texttt{PreIsotonicProfileLikelihood})
exceeds $\Delta_{\max}$ (\texttt{--max\_pl\_jump}, default $1\,\sigma$):
\begin{equation}
  \text{spiked} \;\Leftrightarrow\;
  \max_k\!\left[Z_{\mathrm{pre}}(T_{k+1})-Z_{\mathrm{pre}}(T_k)\right]>\Delta_{\max}.
  \label{eq:spike_criterion}
\end{equation}
The best cut maximizes the reference significance over non-spiked cuts:
\begin{equation}
  (N^*,N_{\mathrm{op}}^*,N_{\mathrm{adj}}^*) =
  \arg\max_{\text{not spiked}} Z_{\mathrm{ref}}(T_{\max}).
  \label{eq:best_cut}
\end{equation}
Spiked cuts are saved to \texttt{\{config\}\_\{name\}\_highest\_spiked\_HEP.pkl}
for diagnostics.

% ===========================================================================
\section{Sensitivity Analysis}
\label{sec:sensitivity}

\subsection{Overview}
\label{sec:sens_overview}

The Sensitivity analysis maps how well the detector can distinguish between
the solar ($\Delta m^2_\odot$) and reactor ($\Delta m^2_{\mathrm{react}}$)
best-fit oscillation hypotheses.  Rather than computing a single discovery
significance, it evaluates a $\chi^2$ surface over the oscillation-parameter grid
$(\Delta m^2,\,\sin^2\theta_{13},\,\sin^2\theta_{12})$.

The analysis is the most complex of the three:
\begin{itemize}[leftmargin=*]
  \item Histograms are \textbf{two-dimensional} (energy $\times$ azimuth $\cos\eta$).
  \item Two normalization nuisances are floated: $\Apred$ (signal) and $\Abkg$ (background).
  \item The nuisance equations have \textbf{no analytic solution}
    (unlike HEP's quadratic for $\bhat$); full 2D numerical minimization is used.
  \item The goal is a $\chi^2$ \emph{map}, not a single significance value.
\end{itemize}

\subsection{Two-Dimensional Template Construction}
\label{sec:sens_templates}

Signal and background are represented as 2D arrays with axes
energy ($j=1,\ldots,J$) and azimuth ($i=1,\ldots,I$).

\medskip
\noindent\textbf{Signal template.}
For oscillation parameters $\vec{\theta} = (\Delta m^2,\sin^2\theta_{13},\sin^2\theta_{12})$,
the signal template at exposure $\mathcal{E}$ is:
\begin{equation}
  \Tsig_{ij}(\vec{\theta}) = \mathcal{E}\cdot M_{\mathrm{det}}
    \cdot \bigl[P(\vec{\theta})\,H\bigr]_{ij},
  \label{eq:signal_template}
\end{equation}
where $P(\vec{\theta})$ is the oscillation-probability matrix
(azimuth $\times$ neutrino-flavor) and $H$ is the detector energy-response
matrix (neutrino energy $\times$ reconstructed energy), implemented in
\texttt{14SensitivitySignalTemplate.py}.

\medskip
\noindent\textbf{Background template.}
The background template $\Tbkg_{ij}$ is independent of oscillation parameters
and is constructed from detector simulations in
\texttt{14SensitivityBackgroundTemplate.py}.

\subsection{Asimov Grid Construction}
\label{sec:sens_asimov}

For each oscillation point $\vec{\theta}_k$ in the scan grid,
the Asimov (fake) observed dataset is:
\begin{equation}
  o_{ij}(\vec{\theta}_k) = \Tsig_{ij}(\vec{\theta}_k) + \Tbkg_{ij}.
  \label{eq:sens_asimov}
\end{equation}
This is the same median-experiment construction as the HEP profile-likelihood
(Section~\ref{sec:pl}): no statistical fluctuations, purely the expected counts.

Two \emph{reference} templates are fixed at the solar and reactor best-fit points:
\begin{equation}
  p^{\mathrm{solar}}_{ij} = \Tsig_{ij}(\vec{\theta}_{\odot}),
  \qquad
  p^{\mathrm{react}}_{ij} = \Tsig_{ij}(\vec{\theta}_{\mathrm{react}}).
  \label{eq:reference_templates}
\end{equation}

\subsection{Objective Function: Baker-Cousins Poisson Deviance}
\label{sec:sens_objective}

The fit minimizes a Baker-Cousins Poisson deviance \cite{BakerCousins1984}
with two free normalization nuisances:
\begin{equation}
  \chi^2(\Apred,\Abkg) = 2\sum_{i,j}\Delta\ell_{ij}
    + \left(\frac{\Apred}{\spred}\right)^{\!2}
    + \left(\frac{\Abkg}{\sbkg}\right)^{\!2},
  \label{eq:sens_chisq}
\end{equation}
where the expected model is:
\begin{equation}
  e_{ij} = (1+\Abkg)\,\Tbkg_{ij} + (1+\Apred)\,p_{ij},
  \label{eq:sens_expected}
\end{equation}
and the per-bin Poisson deviance term is:
\begin{equation}
  \Delta\ell_{ij} =
  \begin{cases}
    e_{ij} - o_{ij} + o_{ij}\ln(o_{ij}/e_{ij}) & o_{ij}>0,\; e_{ij}>0,\\
    e_{ij} & o_{ij}=0,\; e_{ij}>0,\\
    0 & e_{ij}=0 \;\text{(or masked)}.
  \end{cases}
  \label{eq:sens_deviance}
\end{equation}

Each term $\Delta\ell_{ij} \ge 0$ by the Gibbs inequality; the sum is zero
if and only if $e_{ij} = o_{ij}$ for all bins.
This is the direct 2D generalization of the per-bin LLR in the HEP
profile-likelihood (Eq.~\eqref{eq:q0_expanded}), with two nuisances
instead of one.

\subsection{Nuisance Parameters and Constraints}
\label{sec:sens_nuisances}

The penalties $({\Apred}/{\spred})^2 + ({\Abkg}/{\sbkg})^2$ are
Gaussian constraints on the normalization shifts, identical in form to
the HEP penalty $((\bhat-1)/\srel)^2$.
Typical values: $\spred = 4\%$ (signal flux uncertainty),
$\sbkg = 2\%$ (background uncertainty).

\medskip
\noindent\textbf{Why no analytic solution exists.}
In HEP, the background nuisance $\beta$ is global: $e_i = \beta b_i$ (under $H_0$),
so $\sum_i b_i$ and $\sum_i n_i$ are the only sums needed and the stationarity
condition is a quadratic in $\beta$.

In Sensitivity, $e_{ij} = (1+\Abkg)B_{ij} + (1+\Apred)P_{ij}$ contains both
$B$ and $P$ in every bin.  The stationarity conditions are:
\begin{align}
  \sum_{ij} P_{ij}\!\left(1-\frac{o_{ij}}{e_{ij}}\right)
    &= -\frac{\Apred}{\spred^2},
  \label{eq:sens_stat_pred}\\
  \sum_{ij} B_{ij}\!\left(1-\frac{o_{ij}}{e_{ij}}\right)
    &= -\frac{\Abkg}{\sbkg^2},
  \label{eq:sens_stat_bkg}
\end{align}
which are jointly non-linear in $(\Apred,\Abkg)$ because $e_{ij}$ in the
denominator depends on both.  Unlike the HEP case, there is no factorization
into a scalar equation.

\subsection{Barlow-Beeston Mask for Templates}
\label{sec:sens_bb}

Bins where the background template is zero ($\Tbkg_{ij}=0$) have no MC
support and are excluded from the fit:
\begin{equation}
  \mathcal{M}_{ij} = \mathbf{1}\!\left[\Tbkg_{ij} > 0\right],
  \label{eq:sens_bb_mask}
\end{equation}
passed as \texttt{bb\_mask} to \texttt{Sensitivity\_Fitter}.
This prevents spurious large deviance contributions from signal-only bins
with unreliable background estimates, following the same spirit as
HEP's Barlow-Beeston mask (Section~\ref{sec:bb_mask}).

\subsection{Minimization: scipy L-BFGS-B}
\label{sec:sens_minimization}

The objective (Eq.~\eqref{eq:sens_chisq}) is minimized numerically over
$(\Apred,\Abkg)$ jointly using the L-BFGS-B algorithm
\cite{ByrdEtAl1995,ZhuEtAl1997}:
\begin{equation}
  (\hat{A}_{\mathrm{pred}},\hat{A}_{\mathrm{bkg}})
  = \arg\min_{\Apred,\Abkg}\chi^2(\Apred,\Abkg),
  \label{eq:sens_minimize}
\end{equation}
with box constraints:
\begin{equation}
  |\Apred| \le 10\,\spred,
  \qquad
  |\Abkg| \le 10\,\sbkg.
  \label{eq:sens_bounds}
\end{equation}
L-BFGS-B uses gradient information (finite-difference approximation)
and is well-suited to smooth, convex objectives with box constraints.
It replaces a previous 2D Minuit minimization, requiring fewer function
evaluations for the same convergence tolerance.
Implemented in \texttt{lib/lib\_root.py:Sensitivity\_Fitter.Fit}.

\subsection{Oscillation Grid Scan and Best-Cut Score}
\label{sec:sens_scan}

For each analysis cut $(N_{\mathrm{hits}},N_{\mathrm{ophits}},N_{\mathrm{adjcl}})$
and each oscillation point $\vec{\theta}_k$:
\begin{enumerate}[leftmargin=*]
  \item Construct the Asimov dataset $o_{ij}(\vec{\theta}_k)$.
  \item Minimize $\chi^2$ against the solar reference template $p^{\mathrm{solar}}$
    to obtain $\chi^2_{\odot}(\vec{\theta}_k)$.
  \item Minimize $\chi^2$ against the reactor reference template $p^{\mathrm{react}}$
    to obtain $\chi^2_{\mathrm{react}}(\vec{\theta}_k)$.
\end{enumerate}

The \textbf{cut quality score} is the average $\chi^2$ when fitting with the
wrong hypothesis (larger = better discrimination):
\begin{equation}
  \mathrm{Score}(\mathrm{cut}) =
    \tfrac{1}{2}\!\left[
      \chi^2_{\odot}(\vec{\theta}_{\mathrm{react}})
      + \chi^2_{\mathrm{react}}(\vec{\theta}_{\odot})
    \right].
  \label{eq:sens_score}
\end{equation}
The optimal cut maximizes this score.  The resulting $\chi^2$ surface over
the oscillation grid is used to derive sensitivity contours.

% ===========================================================================
\section{Comparison Across Analyses}
\label{sec:comparison}

Table~\ref{tab:comparison} summarizes the statistical model for each analysis.

\begin{table}[ht]
\centering
\caption{Statistical model comparison: Day-Night, HEP, and Sensitivity analyses.}
\label{tab:comparison}
\begin{tabular}{llll}
\toprule
Feature & Day-Night & HEP & Sensitivity \\
\midrule
Observable & 1D asymmetry spectrum & 1D energy spectrum & 2D (energy $\times$ azimuth) \\
Signal & $s_i = \mathcal{E}\xi(r_i^{\rm night}-r_i^{\rm day})$ & $s_i = \mathcal{E}r_i^{\rm sig}$ & $\Tsig_{ij}(\vec{\theta})$ \\
Goal & Detect day-night excess & Discover HEP flux & Map $(\Delta m^2,\sin^2\theta)$ contour \\
Significance type & Gaussian only & Gaussian, Asimov, PL & $\chi^2$ surface \\
Nuisances & None & $1\times\beta$ (background) & $\Apred + \Abkg$ \\
Nuisance solution & --- & Closed-form quadratic & scipy L-BFGS-B (2D) \\
Penalty form & --- & $[(\bhat-1)/\srel]^2$ & $(\Apred/\spred)^2+(\Abkg/\sbkg)^2$ \\
Per-bin deviance & $s_i/\sqrt{b_i}$ & $n_i\ln(n_i/\bhat b_i)-(n_i-\bhat b_i)$ & $e_{ij}-o_{ij}+o_{ij}\ln(o_{ij}/e_{ij})$ \\
Barlow-Beeston mask & No & $N_i^{\rm MC}\ge N_{\rm MC}^{\min}$ & $\Tbkg_{ij}>0$ \\
Earth density band & $\xi\in\{1\pm\varepsilon,1\}$ & --- & --- \\
PL post-processing & --- & Gauss($\sigma=6$) + PAVA & --- \\
Spike filter & --- & $\Delta_{\max}$ on pre-PAVA & --- \\
Output & $Z$ vs exposure & $Z$ vs exposure & $\chi^2$ map over grid \\
\bottomrule
\end{tabular}
\end{table}

\medskip
\noindent\textbf{Narrative progression.}
The three analyses represent increasing statistical complexity.
The Day-Night analysis is the simplest possible approach: a signal defined as
a spectral difference, evaluated with a Gaussian approximation.
No nuisance parameters are profiled; the method is robust but conservative.

The HEP analysis introduces the full profile-likelihood framework.
The key insight is that a single global background nuisance $\beta$,
fully correlated across bins, leads to a quadratic stationarity condition
with an analytic solution.  This makes the PL computation fast and exact.
Post-processing (Gaussian smoothing + PAVA) is needed to recover a smooth,
monotone significance curve from a discrete numerical profile.

The Sensitivity analysis takes the same Baker-Cousins Poisson deviance
as the per-bin LLR in HEP, but extends it to 2D and adds a second nuisance.
The two nuisances are coupled per-bin, breaking the factorization that
enabled the HEP quadratic.  Full 2D minimization is required.
The result is not a significance but a $\chi^2$ map used to derive
sensitivity contours and select optimal analysis cuts.

% ===========================================================================
\section{Summary of Significance Outputs}
\label{sec:summary}

Table~\ref{tab:methods} summarizes all significance computation paths.

\begin{table}[ht]
\centering
\caption{Summary of significance output columns across analyses.}
\label{tab:methods}
\begin{tabular}{llllll}
\toprule
Analysis & Label & Histogram & Rebinning & Method & Monotone? \\
\midrule
\multicolumn{6}{l}{\textit{Day-Night}} \\
 & \texttt{RawGaussian}       & Raw     & None & Gaussian         & No \\
 & \texttt{RawErrorGaussian}  & Raw     & None & Gaussian+Poisson & No \\
 & \texttt{Gaussian}          & Smoothed & None & Gaussian        & No \\
 & \texttt{ErrorGaussian}     & Smoothed & None & Gaussian+Poisson & No \\
\midrule
\multicolumn{6}{l}{\textit{HEP}} \\
 & \texttt{RawGaussianNoRebin}    & Raw    & None    & Gaussian & No \\
 & \texttt{RawAsimovNoRebin}      & Raw    & None    & Asimov   & No \\
 & \texttt{RawGaussian}           & Raw    & Adaptive & Gaussian & Yes \\
 & \texttt{RawAsimov}             & Raw    & Adaptive & Asimov   & Yes \\
 & \texttt{GaussianNoRebin}       & Smooth & None    & Gaussian & No \\
 & \texttt{AsimovNoRebin}         & Smooth & None    & Asimov   & No \\
 & \texttt{Gaussian}              & Smooth & Adaptive & Gaussian & Yes \\
 & \texttt{Asimov}                & Smooth & Adaptive & Asimov   & Yes \\
 & \texttt{RawProfileLikelihood}  & Raw    & None    & PL       & Yes (post) \\
 & \texttt{ProfileLikelihood}     & Smooth & None    & PL       & Yes (post) \\
\midrule
\multicolumn{6}{l}{\textit{Sensitivity}} \\
 & \texttt{chi2\_solar/react} & 2D smooth & None & Baker-Cousins & --- \\
\bottomrule
\end{tabular}
\end{table}

% ===========================================================================
\section{Workflow Flags}
\label{sec:flags}

\begin{table}[ht]
\centering
\caption{Orchestrator flags in \texttt{10SensitivityAnalysis.py}.}
\label{tab:flags}
\begin{tabular}{llp{8cm}}
\toprule
Flag & Default & Effect when disabled \\
\midrule
\texttt{--computation}     & True & Skip all computation; run only plot macros \\
\texttt{--significance}    & True & Skip \texttt{12DayNight.py}, \texttt{13HEP.py}, \texttt{14Sensitivity*.py} \\
\texttt{--fiducialization} & True & Skip \texttt{0XFiducializeSignal.py}, \texttt{0YBestFiducial.py} \\
\texttt{--rebin}           & True & Skip \texttt{11AnalysisSignal.py} adaptive rebinning \\
\bottomrule
\end{tabular}
\end{table}

% ===========================================================================
\section*{References}
\addcontentsline{toc}{section}{References}
\begin{thebibliography}{9}

\bibitem{Cowan2010}
G.~Cowan, K.~Cranmer, E.~Gross, O.~Vitells,
\textit{Asymptotic formulae for likelihood-based tests of new physics},
Eur.\ Phys.\ J.\ C \textbf{71} (2011) 1554.
\url{https://arxiv.org/abs/1007.1727}

\bibitem{BakerCousins1984}
S.~Baker, R.D.~Cousins,
\textit{Clarification of the use of chi-square and likelihood functions in fits to histograms},
Nucl.\ Instrum.\ Meth.\ \textbf{221} (1984) 437--442.
\url{https://doi.org/10.1016/0167-5087(84)90016-4}

\bibitem{BarlowBeeston1993}
R.~Barlow, C.~Beeston,
\textit{Fitting using finite Monte Carlo samples},
Comput.\ Phys.\ Commun.\ \textbf{77} (1993) 219--228.
\url{https://doi.org/10.1016/0010-4655(93)90005-W}

\bibitem{RootHistFactory}
ROOT Collaboration,
\textit{HistFactory: A tool for creating statistical models for use with RooFit and RooStats},
CERN-OPEN-2012-016.
\url{https://root.cern.ch/doc/master/classRooStats_1_1HistFactory_1_1Measurement.html}

\bibitem{Robertson1988}
T.~Robertson, F.T.~Wright, R.L.~Dykstra,
\textit{Order Restricted Statistical Inference},
Wiley, 1988.

\bibitem{ByrdEtAl1995}
R.H.~Byrd, P.~Lu, J.~Nocedal, C.~Zhu,
\textit{A limited memory algorithm for bound constrained optimization},
SIAM J.\ Sci.\ Comput.\ \textbf{16} (1995) 1190--1208.

\bibitem{ZhuEtAl1997}
C.~Zhu, R.H.~Byrd, P.~Lu, J.~Nocedal,
\textit{Algorithm 778: L-BFGS-B: Fortran subroutines for large-scale bound-constrained optimization},
ACM Trans.\ Math.\ Softw.\ \textbf{23} (1997) 550--560.

\end{thebibliography}

\end{document}
